% Generated by scripts/assemble.py; corrections belong in source/corrections.json.
\examhead{133}{华中科技大学 2025 年数学分析试卷}{127}
% p128-b01; source PDF page 128
\par\Needspace{5\baselineskip}
\blockanchor{p128-b01}
\majorquestion{一}计算题. 每题 10 分, 共 40 分.\par
% p128-b02; source PDF page 128
\blockanchor{p128-b02}
\question{1}求极限 $\lim_{x\to 0}\frac{(1 + x^2)^3 - \cos^2x}{\tan x^2}.$\par
% p128-b03; source PDF page 128
\blockanchor{p128-b03}
\question{2}讨论级数 $\sum_{n=1}^{\infty} \int_{0}^{\frac{\pi}{n}} \frac{\sin x}{1 + x^{2024}} \, \mathrm{d}x$ 的敛散性.\par
% p128-b04; source PDF page 128
\blockanchor{p128-b04}
\question{3}设 $\beta > 0, \alpha > 1$ 为常数, 计算无穷积分 $I = \int_{0}^{+\infty} \frac{\mathrm{d}x}{1 + \beta x^{\alpha}}$.\par
% p128-b05; source PDF page 128
\blockanchor{p128-b05}
\question{4}设 $\alpha > 0$, 求极限 $\lim_{x \to +\infty} \left[ (x + 1)^{\alpha} - x^{\alpha} - \frac{x + 1}{x + e^{\sin x}} \right]$.\par
% p128-b06; source PDF page 128
\par\Needspace{5\baselineskip}
\blockanchor{p128-b06}
\majorquestion{二}解答题. 每题 15 分, 共 60 分.\par
% p128-b07; source PDF page 128
\blockanchor{p128-b07}
\question{5}求极限 $\lim_{n\to \infty}\sum_{k = 1}^{n}\frac{n^2}{n^3 + k^3 + k^2}.$\par
% p128-b08; source PDF page 128
\blockanchor{p128-b08}
\question{6}求极限 $\lim_{x\to 0^{+}}\frac{e^{x} - 1 - x}{\sqrt{1 - x} - \cos\sqrt{x}}.$\par
% p128-b09; source PDF page 128
\par\Needspace{6\baselineskip}
\blockanchor{p128-b09}
\question{7}求积分\par
\nopagebreak[4]
% p128-b10; source PDF page 128
\blockanchor{p128-b10}
\displayformula{I = \oint_ {C} (x + y) ^ {2} \mathrm{d} x - (x ^ {2} + y ^ {2}) \mathrm{d} y.}
% p128-b11; source PDF page 128
\blockanchor{p128-b11}
其中 C 是从点  $A(1,1)$  到点  $B(2,5)$  再到点  $C(3,2)$  最后到点  $A(1,1)$  的三角形边界.\par
% p128-b12; source PDF page 128
\par\Needspace{6\baselineskip}
\blockanchor{p128-b12}
\question{8}设 $a, b, c$ 为正常数, $S$ 为单位球面 $x^{2} + y^{2} + z^{2} = 1$ 的外侧, 计算第二类曲面积分\par
\nopagebreak[4]
% p128-b13; source PDF page 128
\blockanchor{p128-b13}
\displayformula{I = \iint_ {S} \frac {x \mathrm{d} y \mathrm{d} z + y \mathrm{d} z \mathrm{d} x + z \mathrm{d} x \mathrm{d} y}{(a x ^ {2} + b y ^ {2} + c z ^ {2}) ^ {\frac {3}{2}}}.}
% p128-b14; source PDF page 128
\par\Needspace{5\baselineskip}
\blockanchor{p128-b14}
\majorquestion{三}证明题. 前两题各 10 分, 后两题各 15 分.\par
% p128-b15; source PDF page 128
\blockanchor{p128-b15}
\question{9}用 $\varepsilon - N$ 语言证明极限 $\lim_{n\to \infty}\frac{1}{\sqrt[n]{(2n)!}} = 0.$\par
% p128-b16; source PDF page 128
\par\Needspace{6\baselineskip}
\blockanchor{p128-b16}
\question{10}设\par
\nopagebreak[4]
% p128-b17; source PDF page 128
\blockanchor{p128-b17}
\displayformula{\begin{array}{c} x _ {1} = 1, x _ {2} = \sqrt {\frac {1}{2} + 1}, x _ {3} = \sqrt {\frac {1}{3} + \sqrt {\frac {1}{2} + 1}}, x _ {4} = \sqrt {\frac {1}{4} + \sqrt {\frac {1}{3} + \sqrt {\frac {1}{2} + 1}}}, \\ \dots , x _ {n + 1} = \sqrt {\frac {1}{n + 1} + x _ {n}}, \dots . \end{array}}
% p128-b18; source PDF page 128
\blockanchor{p128-b18}
证明数列 $\{x_{n}\}$ 收敛并求出极限值.\par
% p128-b19; source PDF page 128
\blockanchor{p128-b19}
\question{11}设函数 $f(x) \in C(0, +\infty)$，满足 $x = f(x)5^{f(x)}$. 证明：\par
% p128-b20; source PDF page 128
\blockanchor{p128-b20}
(1) $f(x)$ 单调递增.\par
% p128-b21; source PDF page 128
\blockanchor{p128-b21}
(2) $\lim_{x\to +\infty}f(x) = +\infty .$\par
% p128-b22; source PDF page 128
\blockanchor{p128-b22}
(3) $\lim_{x\to +\infty}\frac{f(x)}{\ln x} = \frac{1}{\ln 5}.$\par
% p128-b23; source PDF page 128
\par\Needspace{6\baselineskip}
\blockanchor{p128-b23}
\question{12}对常数 $a > 0$, 记平面区域 $D: x^2 + y^2 \leq a^2$ 的边界为 $\partial D$. 设二元函数 $f(x, y)$ 在 $D$ 上有连续偏导数, 且 $f(x, y) = 0$, $(x, y) \in \partial D$. 证明:\par
\nopagebreak[4]
% p128-b24; source PDF page 128
\blockanchor{p128-b24}
\displayformula{\left| \iint_ {D} f (x, y) \mathrm{d} x \mathrm{d} y \right| \leq \frac {\pi a ^ {2}}{3} \max _ {(x, y) \in D} \sqrt {\left(\frac {\partial f}{\partial x}\right) ^ {2} + \left(\frac {\partial f}{\partial y}\right) ^ {2}}.}
